% ID: SEL-D0014-BD-E07
% ESTADO: cotejada. Segunda pasada visual completa y borrador compilado revisados.
% Fuente: Selectividad/matematicas_aplicadas/matematicas_aplicadas_P_B.pdf
% SHA-256: 51601c4f5a6c1dfcd860f621a659f8b9c09f53ff4cd31a1fa1e827dbd8f06159
% Página: 2; bloque D; ejercicio original 7.

El precio de venta al público del kilogramo de frambuesas sigue una ley
Normal de media desconocida y varianza 9. En una localidad se eligen 10
comercios de manera aleatoria, obteniéndose los siguientes precios en euros:
\[
12.3\quad 10\quad 9.1\quad 11\quad 10.5\quad 11.8\quad 9.9\quad
11.5\quad 10.9\quad 13
\]
\begin{enumerate}[label=\alph*)]
  \item \textbf{(0.5 puntos)} \textquestiondown{}Qué distribución siguen
  las medias de las muestras de tamaño 10?
  \item \textbf{(1 punto)} Con los datos obtenidos de la muestra, determine
  un intervalo de confianza al 97\% para el precio medio del kilogramo de frambuesas.
  \item \textbf{(1 punto)} Con el mismo nivel de confianza, calcule el tamaño
  mínimo que debe tener una muestra para que el error cometido al estimar
  el precio medio del kilogramo de frambuesas sea menor a 1.5 euros.
\end{enumerate}
